% =======================================================
% 1. INTRODUZIONE
% =======================================================
\section{Introduzione}

\subsection{Scopo del documento}
Il presente documento, denominato Piano di Qualifica\textsubscript{G}, ha l'obiettivo di definire le strategie, le procedure e gli standard adottati dal gruppo Synergo Unit per garantire la qualità sia dei processi interni che del prodotto finale. \\ \\
Esso specifica le metriche utilizzate, i criteri di accettazione e le attività di verifica e validazione necessarie per assicurare che il software sviluppato soddisfi i requisiti concordati con la proponente Zucchetti\textsubscript{G}.

\subsection{Scopo del prodotto}
Il prodotto, denominato \textbf{Second Brain\textsubscript{G}}, si pone l'obiettivo di rivoluzionare la gestione della conoscenza personale attraverso l'integrazione di Modelli di Linguaggio di Grande Scala (LLM\textsubscript{G}). \\
Il sistema permetterà di:
\begin{itemize}
	\item Organizzare e sintetizzare informazioni da note scritte in formato Markdown\textsubscript{G}.
	\item Estrarre concetti chiave, generare riassunti e traduzioni automatiche.
	\item Fornire funzionalità di supporto alla scrittura tramite interazione con modelli LLM.
	\item Analizzare le note prodotte dall'utente tramite un processo di critica automatico (Sei Cappelli per Pensare\textsubscript{G}, E. De Bono).
\end{itemize}
L'obiettivo finale è supportare, tra le altre cose, il 
\textit{Distant Writing\textsubscript{G}}, permettendo all'utente di 
focalizzarsi sull'elaborazione concettuale delegando all'IA i compiti di 
trasformazione e recupero delle informazioni.
Le metriche e le attività di controllo definite nel presente documento 
sono inoltre utilizzate come strumenti di supporto al monitoraggio dei rischi individuati nel \link{https://synergo-unit-unipd.github.io/Second-Brain/docs/RTB/doc_gestionali/doc_esterni/piano_di_progetto/piano_di_progetto.pdf}{Piano di Progetto}, 
con particolare riferimento alla stabilità dei requisiti, 
all'avanzamento delle attività, alla qualità degli artefatti e 
all'efficacia del processo di verifica.

\subsection{Glossario}
Al fine di evitare ambiguità terminologiche, i termini tecnici, gli acronimi e le parole che potrebbero avere interpretazioni multiple sono definiti nel documento \link{https://synergo-unit-unipd.github.io/Second-Brain/docs/RTB/doc_gestionali/doc_interni/glossario/glossario.pdf}{Glossario}. Tali termini sono contrassegnati nel testo con il pedice \textsubscript{G}.

\subsection{Riferimenti}
\subsubsection{Normativi}
\begin{itemize}
	\item \link{https://www.math.unipd.it/~tullio/IS-1/2025/Progetto/C6.pdf}{Capitolato C6 - Progetto Second Brain}
	\item \link{https://www.math.unipd.it/~tullio/IS-1/2025/Dispense/PD1.pdf}{Regolamento progetto didattico}
	\item \link{https://synergo-unit-unipd.github.io/Second-Brain/docs/RTB/doc_gestionali/doc_interni/norme_di_progetto/norme_di_progetto.pdf}{Norme di Progetto\textsubscript{G}}
	\item
    \link{https://synergo-unit-unipd.github.io/Second-Brain/docs/rif_info_ISO/ISO_12207-1995.pdf}{Standard ISO/IEC 12207:1995} --
    \textit{Information Technology -- Software Life Cycle Processes};
	\item
    \link{https://synergo-unit-unipd.github.io/Second-Brain/docs/rif_info_ISO/ISO-UNI-9001_2015.pdf}{Standard ISO/IEC 9001:2015} --
    \textit{Sistemi di gestione per la qualità -- Fondamenti e vocabolario};
	\item 
	\link{https://www.iso.org/obp/ui/en/\#iso:std:iso-iec:25010:ed-1:v1:en}{Standard ISO/IEC 25010:2011} -- 
	\textit{Systems and software Quality Requirements and Evaluation (SQuaRE) -- System and software quality models};
\end{itemize}
\subsubsection{Informativi}
\begin{itemize}
	\item \link{https://synergo-unit-unipd.github.io/Second-Brain/docs/RTB/doc_tecnici/doc_esterni/analisi_dei_requisiti/analisi_dei_requisiti.pdf}{Analisi dei Requisiti\textsubscript{G}}
	\item \link{https://synergo-unit-unipd.github.io/Second-Brain/docs/RTB/doc_gestionali/doc_esterni/piano_di_progetto/piano_di_progetto.pdf}{Piano di Progetto\textsubscript{G}}
	\item Slide del corso di Ingegneria del Software

\end{itemize}